\section{Experimental Setup}\label{experimental_setup}

\subsection{Data and Splits}

Experiments were conducted using the GQA dataset which is based on the Visual Genome dataset \parencite{hudson2019gqa}.
A subset of 5000 question-answer pairs were used for the study. To prevent information leakage between
splits, examples were partitioned at the image level such that the same image did not occur in more than 
one split. This resulted in 3479 training examples, 753 validation examples, and 768 test examples.

Only questions that were suitable for localised privileged visual context were retained. Questions requiring
broader spatial context, complex relational reasoning or comparison across distant objects were removed. This
was done using GQA metadata which provide information about question semantic types, structural types and functional
programs. Textual spatial cues were also used to refine what questions to use. Lastly, questions were
restricted to a maximum of two grounded objects.

For each retained question, the union of the corresponding grounded object bounding boxes was used to 
construct the privileged teacher crop. Crops were expanded by 12\% to preserve local surrounding context.
Candidate crops were required to occupy between 2\% and 50\% of the original image before padding and no more
than 55\% after padding. A minimum crop side length of 24 pixels was also enforced. No answer-side information
was used during crop construction.

\subsection{Vision-Language Models}

Three VLMs with different architectures and capacities were evaluated:

\begin{itemize}
\item \textbf{SmolVLM-500M-Instruct}, representing a compact generative vision-language model;
\item \textbf{BLIP-VQA-Base} \cite{li2022blip}, representing an encoder-decoder VQA architecture; and
\item \textbf{MiniCPM-V 4.6}, representing a larger modern multimodal generative model.
\end{itemize}

Each architecture was evaluated under three conditions: the original pretrained model, SFT and OPSD.
The pretrained condition provides a zero-shot reference, while SFT and PSD use the same training
split and parameter-efficient adaptation framework.

\subsection{Training Configuration}

Both SFT and OPSD were implemented using LoRA \parencite{hu2021lora} with rank $r=8$ and scaling
parameter $\alpha=16$. Only LoRA parameters were updated, while the underlying pretrained model
weights remained frozen. Models were trained for three epochs using two NVIDIA T4 GPUs from 
Kaggle with DistributedDataParallel. A per-device batch size of 1 and gradient accumulation
over four steps were used, giving an effective batch size of $1 \times 2 \times 4 = 8$.

The SFT learning rate was set to $1\times10^{-4}$, while the OPSD learning rate was set to $5\times10^{-6}$. The privileged teacher remained frozen
while the full-image student LoRA parameters were updated. No supervised cross-entropy anchor was included
in the OPSD objective. The Jensen-Shannon divergence was computed with equal weighting between the student
and the teacher distributions.

For each epoch, model performance was evaluated on the validation set. The checkpoint with the
highest validation BETScore-F1 was selected as the final model for both SFT and OPSD. The held-out 
test set was not used for hyperparameter selection or checkpoint selection.

For each architecture, the following quantities were reported:

\[F1_{\mathrm{Base}}~, F1_{\mathrm{SFT}},~ F1_{\mathrm{OPSD}}.\]

The final quantity directly addresses RQ1, while comparing it across the three architectures
addresses RQ2.

Because SFT and OPSD were evaluated on the same test questions, an additional paired analysis was
performed at the question level. For each example $i$, 

\[F1^{(i)}_{\mathrm{OPSD}},
F1^{(i)}_{\mathrm{SFT}}.\]

The mean paired difference was reported together with a 95\% bootstrap confidence interval using
10000 resamples. The proportions of test examples for which OPSD outperformed SFT, SFT outperformed OPSD or 
the two methods tied were also recorded.